% ============================================================
% 二、问题形式化 (Problem Formalization)
% ============================================================
\section{问题形式化}

\subsection{变量与符号定义}

用严格的数学符号定义问题中涉及的所有变量。建议使用表格整理：

\begin{table}[htbp]
    \centering
    \caption{符号定义表}
    \label{tab:notation}
    \begin{tabular}{c|l}
        \toprule
        符号 & 含义 \\
        \midrule
        $\X \subseteq \R^d$ & 输入/数据空间 \\
        $\Y$                & 输出/标签空间 \\
        $\mathbf{x}_i$      & 第 $i$ 个样本 \\
        $\boldsymbol{\theta}$ & 待优化参数 \\
        $\Loss(\cdot)$      & 损失/目标函数 \\
        \bottomrule
    \end{tabular}
\end{table}

\subsection{假设}

明确列出为简化问题或使推导成立所引入的所有假设，每一条都要有明确的数学表述：

\begin{enumerate}[label=(\textbf{A\arabic*}), leftmargin=*]
    \item \textbf{数据假设}：【例如：样本独立同分布、数据满足某种光滑性条件等】
    \item \textbf{模型假设}：【例如：线性模型、高斯噪声、低秩结构等】
    \item \textbf{优化假设}：【例如：目标函数 $L$-光滑、$\mu$-强凸、约束可行域非空等】
\end{enumerate}

\subsection{目标函数与约束}

写出完整的优化问题或数学模型的数学表达式：

\begin{equation}\label{eq:problem}
    \min_{\mathbf{x} \in \X} \; f(\mathbf{x}) \quad \text{s.t.} \quad g_i(\mathbf{x}) \leq 0,\; i = 1,\ldots,m
\end{equation}

或对于统计推断/学习问题，写出概率模型和推断目标。

\subsection{问题类别}

指出问题所属的数学类别（勾选适用的项目）：

\begin{itemize}
    \item 凸 / 非凸
    \item 连续 / 离散 / 混合整数
    \item 确定 / 随机
    \item 约束 / 无约束
    \item 光滑 / 非光滑
    \item 凸优化 / 组合优化 / 变分推断 / 谱方法 / PDE / 泛函分析
\end{itemize}

这一判断直接影响后续方法的选择与分析路径。
